We first isolate two (minor) ambiguities in the statement as extracted.

\begin{itemize}
\item The expression ``$\prod_{1\le i\le \log n}$'' uses a non-integer upper limit.  In what follows I interpret
\[k(n):=\lfloor \log n\rfloor\qquad(\log=\log_e)\]
(and note that changing to $\lceil\log n\rceil$ or a different logarithm base only perturbs constants).
\item The text claims that taking $n$ to be the product of primes between $\log n$ and $(2+o(1))\log n$ gives an example with $q(n,\log n)\ge (2+o(1))\log n$.  For the literal definition
\[q(n,k):=\min\{\text{prime }p: p\nmid \prod_{i=1}^k (n+i)\},\]
this appears inconsistent (Lemma~457.3 below makes this precise).  A ``minimal correction'' that makes the example valid is to replace the product by $\prod_{0\le i\le k}(n+i)$ (include the factor $n$).  In this writeup I treat the \emph{literal} definition from the file, and separately note what changes under the corrected variant.
\end{itemize}

\noindent\textbf{Literal problem.}

Define for integers $n\ge 2$ and $k\ge 0$
\[P(n,k):=\prod_{i=1}^k (n+i)\quad (\text{empty product }=1),\qquad q(n,k):=\min\{\text{prime }p: p\nmid P(n,k)\}.
\]
Question A: Does there exist $\varepsilon>0$ and infinitely many integers $n$ such that
\[q(n,k(n))\ \ge\ (2+\varepsilon)\log n\ ?\]
Equivalently, for infinitely many $n$, \emph{every} prime $p\le (2+\varepsilon)\log n$ divides $P(n,k(n))$.

Question B (from the ``more generally'' clause): Does there exist $\varepsilon>0$ such that for all sufficiently large $n$,
\[q(n,k(n))\ <\ (1-\varepsilon)(\log n)^2\ ?\]
